\subsection*{\texttt{parallel.async}: host-level concurrency}

Everything above runs inside a kernel. Sometimes you need concurrency at the \textbf{host level}: launch a grid without blocking the CPU.

\begin{lstlisting}[language=croqtile,caption={parallel.async: host-level concurrency}]
parallel.async {px, py} by [grid_m, grid_n] : block {
  // kernel body
}
\end{lstlisting}

\textbf{\texttt{parallel.async}} returns control to the host immediately --- the kernel is enqueued but the host does not wait for completion. This is the Croqtile equivalent of \texttt{launch\_kernel} with a non-default stream.

This is \textbf{host orchestration}, orthogonal to in-kernel control flow. It does not replace \texttt{inthreads.\allowbreak{}async} for thread partitioning or \texttt{if} for runtime predicates --- it decides \emph{when} and \emph{where} a grid runs relative to other grids.
